\scanpage{79}
- \mspage{208} -
\margin{n$^\circ$ 391}
(III)\quad En vertu du \S~8, n$^\circ$ 3, formule (4), on a $h=n/\ell=12$.

(IV)\quad L'ordre de la plus grande racine est $h-1=11$, (\S~8, n$^\circ$ 5, prop. 6).
Or $\varepsilon_1+\varepsilon_4=2\alpha_1+3\alpha_2+4\alpha_3+2\alpha_4$ est d'ordre 11, et c'est donc la plus grande
racine $\widetilde{\alpha}$; elle est orthogonale à toutes les racines simples sauf $\alpha_1$ et on a
$(\widetilde{\alpha}|\alpha_1)=1$.

(V)\quad Les \uline{coracines} sont les vecteurs $\pm2\varepsilon_i^*$, $\pm\varepsilon_i^*\pm\varepsilon_j^*$ $(i\ne j)$, et
$\pm\varepsilon_1^*\pm\varepsilon_2^*\pm\varepsilon_3^*\pm\varepsilon_4^*$. Les angles des coracines étant égaux à ceux des racines
correspondantes, le système de racines $R^*$ ne peut avoir comme schéma de Dynkin
que $F_4$, (puisque le schéma de Coxeter détermine alors le schéma de Dynkin), et $R^*$
est donc isomorphe à $R$, ce que l'on peut vérifier directement en calculant les
images canoniques $\alpha_i^*$ des $\alpha_i$ qui forment un système de racines simples de $R^*$. Les
racines non orthogonales à $u=\varepsilon_1$ sont $\pm u$, $\pm\varepsilon_1\pm\varepsilon_j$, $(j\geq2)$ et
$\frac12(\pm\varepsilon_1\pm\varepsilon_2\pm\varepsilon_3\pm\varepsilon_4)$, donc sont au nombre de 30, et on vérifie aussitôt que
$n(u,v)=\pm2$ pour 14 de ces racines et $n(u,v)=\pm1$ pour les 16 autres. La \uline{forme
bilinéaire canonique} pour $R$ est donc $(x|y)/18$ et la plus petite longueur des
racines est $1/3\sqrt2$; le même calcul que pour le type $B_\ell$ donne $\unsure{\gamma|\chi}=2.3^4$.

(VI)\quad Le réseau $\Gamma_o(R^*)$ est le réseau $L_1(V^*)$ et on a vu, (n$^\circ$ 4), que son
associé $\Gamma(R)$ dans $V=R^4$ est $L_2=\Gamma_o(R)$. Le calcul des poids fondamentaux,
formant la base duale de $(\alpha_i^*)$, donne ici :
\[
\begin{aligned}
\omega_1&=\varepsilon_1+\varepsilon_4=2\alpha_1+3\alpha_2+4\alpha_3+2\alpha_4=\widetilde{\alpha}\\
\omega_2&=\varepsilon_1+\varepsilon_2+2\varepsilon_4=3\alpha_1+6\alpha_2+8\alpha_3+4\alpha_4\\
\omega_3&=\tfrac12(\varepsilon_1+\varepsilon_2+\varepsilon_3+3\varepsilon_4)=2\alpha_1+4\alpha_2+6\alpha_3+3\alpha_4\\
\omega_4&=\varepsilon_4=\alpha_1+2\alpha_2+3\alpha_3+2\alpha_4.
\end{aligned}
\]
\mspage{78}
